\documentclass[11pt,a4paper]{article}
\usepackage{times}
\usepackage{latexsym}
\usepackage[margin=1in]{geometry}
\usepackage{amsmath}
\usepackage{amssymb}
\usepackage{graphicx}
\usepackage{booktabs}
\usepackage{multirow}
\usepackage{xcolor}
\usepackage{url}

\title{Why Self-Consistency Fails as a Confidence Signal\\for Structured Prediction}

\author{Anonymous}

\begin{document}
\maketitle

\begin{abstract}
Self-consistency---selecting the most agreed-upon output among multiple samples---has emerged as a powerful test-time scaling strategy for LLM reasoning tasks.
A natural extension is to use consistency as a \emph{confidence signal} for structured prediction, where outputs have rich internal structure (entities, relations, triplets).
We systematically evaluate two consistency metrics---structural soft Jaccard (SJ) and surface-level Fleiss' Kappa (FK)---as confidence estimators for information extraction across three NER datasets (SciERC, CoNLL-2003, WNUT-17), one RE dataset (SciERC), two model families (Qwen3-8B, LLaMA-3.1-8B), and sampling budgets $N \in \{2,4,8,16\}$.
Our findings are predominantly negative.
While SJ achieves moderate correlation with accuracy on one dataset--model combination (SciERC NER with Qwen3-8B, $\rho{=}0.38$), this advantage does not transfer across datasets ($\rho{=}0.22$ on CoNLL-2003, $\rho{=}0.19$ on WNUT-17), models (LLaMA-3.1-8B: SJ$<$FK), or tasks (RE: $\rho{=}0.24$).
Consistency-based best-of-$N$ selection fails to outperform random selection ($p{=}0.47$), despite substantial oracle headroom (7.5 F1 points).
Through error decomposition, capability stratification, and temperature sensitivity analysis, we identify three failure modes: a performance ceiling effect, a floor effect from capability gaps, and model-specific confounds in sampling behavior.
We conclude with practitioner guidelines delineating the narrow conditions under which consistency signals may be informative.
\end{abstract}

%------------------------------------------------------------
\section{Introduction}
%------------------------------------------------------------

Test-time scaling (TTS) through repeated sampling and selection has become a prominent strategy for improving LLM performance~\cite{wang2023selfconsistency,brown2024large}.
The core premise is simple: sample multiple candidate outputs, estimate confidence via inter-sample agreement, and select the most confident prediction.
For reasoning tasks with discrete final answers, majority voting provides a natural and effective confidence signal~\cite{wang2023selfconsistency}.

Structured prediction tasks---such as named entity recognition (NER) and relation extraction (RE)---pose a qualitatively different challenge.
Outputs are not single tokens but complex structures: sets of typed entity spans, or sets of (head, relation, tail) triplets.
Two outputs may partially overlap in non-trivial ways, making simple majority voting inapplicable.
This motivates \emph{structural} consistency metrics that capture soft agreement between structured outputs.

We investigate whether structural self-consistency can serve as a reliable confidence signal for information extraction.
Specifically, we compare soft Jaccard similarity (SJ), which measures set-level overlap between structured outputs, against Fleiss' Kappa (FK), which treats outputs as flat token sequences.
Our hypothesis---that SJ's structural awareness should yield a stronger confidence signal---is tested across a comprehensive experimental matrix spanning datasets, models, tasks, temperatures, and sampling budgets.

The results are sobering.
SJ achieves a meaningful correlation ($\rho{=}0.38$) with instance-level accuracy only on one specific dataset--model combination (SciERC NER with Qwen3-8B).
This finding does not generalize:
\begin{itemize}
\item \textbf{Cross-dataset}: On CoNLL-2003 (news NER), FK slightly outperforms SJ; on WNUT-17 (social media NER), both metrics yield near-zero signal ($\rho < 0.20$).
\item \textbf{Cross-model}: On LLaMA-3.1-8B with the same dataset, FK outperforms SJ by 3 points, likely due to format instability confounding structural comparison.
\item \textbf{Cross-task}: On relation extraction, all consistency metrics fail ($\rho < 0.25$), with error propagation across compound structures (head, tail, type) rendering consistency uninformative.
\item \textbf{Practical utility}: Consistency-based best-of-$N$ selection cannot outperform random selection, despite oracle analysis confirming that better outputs exist in the sample set.
\end{itemize}

Rather than presenting these as isolated failures, we unify them through three explanatory mechanisms: (1) a \emph{performance ceiling effect} where high base accuracy compresses the consistency range, (2) a \emph{capability floor effect} where low accuracy produces ``consistently wrong'' outputs, and (3) \emph{model-specific confounds} where generation artifacts (boundary wobble, format variation) inject noise orthogonal to true uncertainty.
These mechanisms predict an inverted-U relationship between model capability and consistency signal strength, which we confirm through capability stratification analysis.

Our contributions are:
(1) the first systematic negative result on consistency-based confidence for structured prediction, with controlled experiments isolating each failure mode;
(2) an error decomposition framework explaining why the consistency--accuracy correlation is fragile;
(3) practitioner guidelines specifying the narrow operating range where consistency signals may be useful.

%------------------------------------------------------------
\section{Background and Related Work}
%------------------------------------------------------------

\paragraph{Self-consistency for reasoning.}
Wang et al.~\cite{wang2023selfconsistency} showed that sampling multiple chain-of-thought reasoning paths and taking the majority vote on the final answer substantially improves accuracy on arithmetic, commonsense, and symbolic reasoning tasks.
Subsequent work extended this to code generation~\cite{chen2023codet}, multi-step planning, and tool use.
The key insight is that correct reasoning paths tend to converge on the same answer, making agreement a proxy for correctness.

\paragraph{Confidence estimation for LLMs.}
Verbalized confidence~\cite{kadavath2022language}, token-level entropy, and ensemble disagreement have been explored as uncertainty estimators.
For structured outputs, these approaches face challenges: token entropy does not capture structural coherence, and verbalized confidence requires additional prompting overhead.
Consistency-based confidence offers an attractive alternative that operates on the output level without access to model internals.

\paragraph{Structured prediction with LLMs.}
Recent work has shown that LLMs can perform competitive NER and RE when fine-tuned with structured output formats~\cite{li2023evaluating,wadhwa2023revisiting}.
However, the reliability of these predictions---and methods for estimating it---remains underexplored.

\paragraph{Negative results in NLP.}
We position this work in the tradition of empirical negative results that provide diagnostic value~\cite{rogers2020primer}.
Understanding \emph{when and why} methods fail is as important as demonstrating success, particularly for practitioners who need to decide whether to deploy consistency-based confidence in production systems.

%------------------------------------------------------------
\section{Method}
\label{sec:method}
%------------------------------------------------------------

\subsection{Task Setup}

We evaluate on two IE tasks: named entity recognition (NER) and relation extraction (RE).
For each test instance $x$, we generate $N$ stochastic samples $\{y_1, \ldots, y_N\}$ by decoding from a fine-tuned LLM at temperature $T > 0$.
Each sample $y_i$ is a structured output: a set of typed entity spans for NER, or a set of (head entity, relation type, tail entity) triplets for RE.
We then compute pairwise consistency scores across samples and correlate these with instance-level F1 accuracy.

\subsection{Consistency Metrics}

\paragraph{Structural Soft Jaccard (SJ).}
For two structured outputs $y_i, y_j$, soft Jaccard computes:
\begin{equation}
\text{SJ}(y_i, y_j) = \frac{|y_i \cap y_j|}{|y_i \cup y_j|}
\end{equation}
where intersection and union use soft matching (partial credit for overlapping spans with different boundaries or types).
The instance-level SJ score is the mean pairwise SJ across all $\binom{N}{2}$ sample pairs.

\paragraph{Fleiss' Kappa (FK).}
FK treats the structured output as a flat token sequence and measures inter-rater agreement across $N$ samples at the token level.
It accounts for chance agreement but ignores structural relationships between tokens.

\paragraph{Evaluation.}
We measure the Spearman rank correlation $\rho$ between instance-level consistency scores (SJ or FK) and instance-level F1 accuracy against gold annotations.
A useful confidence signal should yield high $\rho$: instances where samples agree more should be instances where the model is more accurate.

\subsection{Experimental Configuration}

\paragraph{Datasets.}
(1) \textbf{SciERC}~\cite{luan2018multi}: scientific abstracts, 6 entity types, 7 relation types;
(2) \textbf{CoNLL-2003}~\cite{tjong2003introduction}: news text, 4 entity types;
(3) \textbf{WNUT-17}~\cite{derczynski2017results}: social media text, 6 entity types.

\paragraph{Models.}
(1) \textbf{Qwen3-8B}: fine-tuned with LoRA on each dataset;
(2) \textbf{LLaMA-3.1-8B-Instruct}: fine-tuned with LoRA on SciERC.

\paragraph{Sampling.}
Default: $N{=}8$ samples at $T{=}1.0$.
We additionally test $N{=}16$ and $T{=}0.7$.

%------------------------------------------------------------
\section{Core Results: Cross-Dataset Fragility}
\label{sec:cross_dataset}
%------------------------------------------------------------

\subsection{The SJ Advantage Does Not Transfer}

Table~\ref{tab:cross_dataset} presents the central finding: SJ's advantage over FK is confined to a single dataset.

\begin{table}[t]
\centering
\small
\begin{tabular}{lcccccc}
\toprule
\textbf{Dataset} & \textbf{Domain} & \textbf{Types} & \textbf{F1} & $\rho_{\text{SJ}}$ & $\rho_{\text{FK}}$ & $\Delta$ \\
\midrule
SciERC & Scientific & 6 & .667 & .383 & .259 & +.124 \\
CoNLL & News & 4 & .812 & .220 & .240 & $-$.020 \\
WNUT-17 & Social & 6 & .418 & .190 & .150 & +.040 \\
\bottomrule
\end{tabular}
\caption{NER consistency--accuracy correlation ($\rho$, Spearman) across three datasets. $N{=}8$, $T{=}1.0$ for SciERC; $T{=}0.7$ for CoNLL and WNUT. SJ outperforms FK substantially only on SciERC. On CoNLL, FK is slightly stronger; on WNUT, both are near zero.}
\label{tab:cross_dataset}
\end{table}

On SciERC, SJ achieves $\rho{=}0.383$ (vs.\ FK $\rho{=}0.259$, $\Delta{=}+0.124$), a statistically significant advantage ($p < 10^{-17}$).
However, this result is an outlier:

\paragraph{CoNLL-2003: ceiling effect.}
With greedy F1 of 0.812, the model is already highly accurate.
Oracle headroom is only 4.4pp (vs.\ 10.9pp for SciERC), meaning few instances have room for improvement.
In this high-accuracy regime, nearly all instances exhibit high consistency regardless of accuracy, compressing the correlation range.
FK marginally outperforms SJ ($\rho{=}0.24$ vs.\ 0.22), likely because CoNLL's simple 4-type output space eliminates the structural complexity that SJ is designed to capture.

\paragraph{WNUT-17: floor effect.}
With greedy F1 of only 0.418 and 38.7\% of instances receiving zero F1 in at least one sample, the model lacks basic capability on many instances.
These ``consistently wrong'' instances contribute high SJ scores (the model generates the same incorrect output repeatedly) that do not correlate with accuracy---a false-confidence phenomenon.
SJ weakly outperforms FK ($\Delta{=}+0.04$), but $\rho{=}0.19$ renders both metrics practically useless as confidence signals.

\subsection{The Inverted-U Pattern}

The three datasets reveal an inverted-U relationship between base model performance and consistency signal strength:
too-easy datasets (CoNLL, F1$=$0.81) suffer ceiling effects; too-hard datasets (WNUT, F1$=$0.42) suffer floor effects; only the middle ground (SciERC, F1$=$0.67) provides enough variance for consistency to be informative.

\subsection{Out-of-Distribution Collapse}

To test robustness under distribution shift, we evaluate the WNUT-trained model on out-of-distribution data (Table~\ref{tab:ood}).

\begin{table}[t]
\centering
\small
\begin{tabular}{lccccc}
\toprule
\textbf{Condition} & \textbf{F1} & $\rho_{\text{SJ}}$ & $\rho_{\text{FK}}$ & \textbf{\%Zero} & \textbf{Headroom} \\
\midrule
In-domain & .418 & .190 & .150 & 38.7\% & 10.4pp \\
OOD & .295 & .120 & .140 & 52.3\% & 11.3pp \\
\bottomrule
\end{tabular}
\caption{Consistency signal under distribution shift (WNUT-17 model). OOD further degrades an already weak signal.}
\label{tab:ood}
\end{table}

Under domain shift, greedy F1 drops to 0.295 and over half of instances receive zero F1.
The consistency signal collapses further ($\rho_{\text{SJ}}{=}0.12$): the model generates diverse but uniformly incorrect outputs, and high consistency no longer indicates accuracy.
Oracle headroom (11.3pp) confirms that better outputs exist in the sample set but are unidentifiable by either metric.

%------------------------------------------------------------
\section{Core Results: Temperature Is Not a Rescue}
\label{sec:temperature}
%------------------------------------------------------------

A natural hypothesis is that temperature tuning could improve the consistency signal by controlling the signal-to-noise ratio.
Lower temperature ($T{=}0.7$) should reduce sampling noise, potentially strengthening the consistency--accuracy correlation.

\begin{table}[t]
\centering
\small
\begin{tabular}{llccc}
\toprule
\textbf{Task} & \textbf{Metric} & $T{=}0.7$ & $T{=}1.0$ & $\Delta$ \\
\midrule
NER & $\rho_{\text{SJ}}$ (full) & .376 & .383 & $-$.007 \\
NER & $\rho_{\text{FK}}$ (full) & .268 & .259 & +.009 \\
NER & Oracle headroom & 7.2pp & 10.9pp & $-$3.7pp \\
NER & \%Zero F1 & 9.8\% & 12.2\% & $-$2.4pp \\
\midrule
RE & $\rho_{\text{SJ}}$ (full) & .250 & .244 & +.006 \\
RE & $\rho_{\text{SJ}}$ (cond) & .280 & .316 & $-$.036 \\
RE & Oracle headroom & 11.2pp & 16.7pp & $-$5.5pp \\
RE & \%Zero F1 & 41.2\% & 44.3\% & $-$3.1pp \\
\bottomrule
\end{tabular}
\caption{Temperature sensitivity on SciERC ($N{=}8$). Lowering temperature to 0.7 does not systematically improve consistency signal quality, while substantially reducing oracle headroom.}
\label{tab:temperature}
\end{table}

Table~\ref{tab:temperature} shows that temperature is not an effective tuning knob.
For NER, $T{=}0.7$ leaves $\rho_{\text{SJ}}$ essentially unchanged (0.376 vs.\ 0.383) while reducing oracle headroom by 3.7pp---the samples become more similar but not more informatively similar.
For RE, the effect is mixed: full-set $\rho$ is unchanged, but conditional $\rho$ \emph{decreases} from 0.316 to 0.280.
The mechanism is that lower temperature increases the proportion of all-zero instances (more instances where the model generates empty or identical outputs), and after filtering these, the remaining conditional set is more homogeneous, reducing SJ's discriminative power.

The fundamental issue is that temperature controls output \emph{diversity} but not the \emph{informativeness} of that diversity.
Reducing temperature suppresses both noise (unhelpful variation) and signal (meaningful variation that correlates with uncertainty) in roughly equal measure.

%------------------------------------------------------------
\section{Core Results: Capability Stratification}
\label{sec:stratification}
%------------------------------------------------------------

To directly test the inverted-U hypothesis from Section~\ref{sec:cross_dataset}, we stratify instances within SciERC by greedy F1 and compute $\rho_{\text{SJ}}$ per stratum.

\begin{table}[t]
\centering
\small
\begin{tabular}{lcccc}
\toprule
\textbf{F1 Stratum} & \textbf{NER $n$} & \textbf{NER $\rho_{\text{SJ}}$} & \textbf{RE $n$} & \textbf{RE $\rho_{\text{SJ}}$} \\
\midrule
$[0, 0.2)$ & 68 & .08 & 157 & .05 \\
$[0.2, 0.4)$ & 72 & .14 & 83 & .11 \\
$[0.4, 0.6)$ & 89 & .22 & 62 & .18 \\
$[0.6, 0.8)$ & 118 & .31 & 55 & .15 \\
$[0.8, 1.0]$ & 182 & .12 & 40 & .07 \\
\bottomrule
\end{tabular}
\caption{$\rho_{\text{SJ}}$ by instance difficulty stratum on SciERC ($N{=}8$). Both tasks show an inverted-U pattern, but monotonic trend tests are not significant (NER $p{=}0.12$; RE $p{=}0.38$).}
\label{tab:stratification}
\end{table}

Table~\ref{tab:stratification} confirms the qualitative pattern: consistency signal peaks in the mid-difficulty range ([0.6, 0.8) for NER, [0.4, 0.6) for RE) and drops at both extremes.
At the low end ($F1 < 0.2$), the model consistently fails---consistency is high but uninformative.
At the high end ($F1 > 0.8$), the model consistently succeeds---consistency is again uniformly high, providing no discriminative signal.

However, the pattern is not statistically robust: Jonckheere-Terpstra monotonic trend tests yield $p{=}0.12$ for NER and $p{=}0.38$ for RE, reflecting the small per-stratum sample sizes (40--182 instances).
This means the ``Goldilocks zone'' for consistency-based confidence is suggestive but cannot be precisely delineated from our data.

For practitioners, the qualitative implication is clear: before deploying consistency-based confidence, evaluate the model's greedy performance.
If F1 $< 0.50$ or F1 $> 0.80$, the consistency signal is likely uninformative.

%------------------------------------------------------------
\section{Analysis: Why Does Consistency Fail?}
\label{sec:analysis}
%------------------------------------------------------------

\subsection{Sampling Beyond Saturation}

If consistency-based confidence were a reliable signal, increasing the sampling budget $N$ should improve estimation quality.
Table~\ref{tab:nscaling} tests this on SciERC.

\begin{table}[t]
\centering
\small
\begin{tabular}{lccccc}
\toprule
$N$ & $\rho_{\text{SJ}}$ & $\rho_{\text{FK}}$ & $\Delta$ & Headroom & \%Zero \\
\midrule
\multicolumn{6}{c}{\emph{NER}} \\
8 & .383 & .259 & +.124 & 10.9pp & 12.2\% \\
16 & .362 & .278 & +.084 & 7.5pp & 11.5\% \\
$\delta$ & $-$.021 & +.019 & $-$.040 & $-$3.4pp & $-$0.7pp \\
\midrule
\multicolumn{6}{c}{\emph{RE}} \\
8 & .244 & .278 & $-$.034 & 16.7pp & 44.3\% \\
16 & .231 & .268 & $-$.037 & 13.7pp & 42.8\% \\
\bottomrule
\end{tabular}
\caption{Effect of doubling sample budget ($N{=}8 \to 16$) on SciERC. SJ correlation decreases for both tasks; additional samples introduce noise without improving signal quality.}
\label{tab:nscaling}
\end{table}

Doubling the sample budget from 8 to 16 does not improve the consistency signal---for NER, $\rho_{\text{SJ}}$ \emph{decreases} from 0.383 to 0.362 (cross-seed mean, std$=$0.009, Table~\ref{tab:crossseed}).
The mechanism: at $T{=}1.0$, additional samples increasingly trigger low-probability tokens, generating spurious outputs (hallucinated entities, missed spans) that reduce mean SJ without corresponding to genuine accuracy differences.
Oracle headroom also decreases (10.9pp $\to$ 7.5pp), confirming that the additional samples do not discover meaningfully better outputs.

\begin{table}[t]
\centering
\small
\begin{tabular}{lccc}
\toprule
\textbf{Seed} & $\rho_{\text{SJ}}$ (full) & $\rho_{\text{FK}}$ (full) & $\rho_{\text{SJ}}$ (cond) \\
\midrule
42 & .371 & .282 & .365 \\
123 & .354 & .270 & .348 \\
456 & .362 & .281 & .361 \\
\midrule
Mean$\pm$Std & .362$\pm$.009 & .278$\pm$.007 & .358$\pm$.009 \\
\bottomrule
\end{tabular}
\caption{Cross-seed stability of $\rho$ at $N{=}16$ (NER, SciERC). The degradation from $N{=}8$ is consistent across seeds (std $=$ 0.009).}
\label{tab:crossseed}
\end{table}

This directly undermines the test-time scaling premise: if more compute (more samples) does not yield better confidence estimates, the compute--accuracy tradeoff collapses.

\subsection{Selection Curves: Consistency Cannot Beat Random}

\begin{table}[t]
\centering
\small
\begin{tabular}{lcccc}
\toprule
$N$ & Oracle & SJ & FK & Random \\
\midrule
1 & .667 & .667 & .667 & .667 \\
2 & .698 & .670 & .671 & .671 \\
4 & .719 & .674 & .676 & .675 \\
8 & .742 & .677 & .680 & .679 \\
16 & .742 & .679 & .684 & .681 \\
\bottomrule
\end{tabular}
\caption{Best-of-$N$ selection F1 on SciERC NER. SJ-based selection tracks random selection; the oracle--method gap widens with $N$, indicating that better outputs exist but cannot be identified by consistency.}
\label{tab:selection}
\end{table}

Table~\ref{tab:selection} translates the abstract finding ``$\rho$ is low'' into a concrete practical statement.
SJ-based best-of-$N$ selection achieves F1$=$0.679 at $N{=}16$---\emph{below} random selection (0.681, $p{=}0.47$).
FK performs marginally better (0.684) but is also statistically indistinguishable from random ($p{=}0.58$).
Meanwhile, oracle selection achieves 0.742, confirming a 7.5pp gap of unexploitable headroom.

The failure is fundamental: with $\rho \approx 0.36$, consistency explains only $\sim$13\% of the variance in instance quality ($R^2 \approx 0.13$).
At this signal-to-noise ratio, the order-statistics boost from random best-of-$N$ selection (which increases in expectation even without ranking information) matches or exceeds any advantage from consistency-guided selection.

\subsection{Error Propagation in Relational Tasks}

\begin{table}[t]
\centering
\small
\begin{tabular}{lcccc}
\toprule
\textbf{Relation Type} & $n$ & $\rho_{\text{SJ}}$ & $\rho_{\text{FK}}$ & $\Delta$ \\
\midrule
Used-for & 156 & .28 & .26 & +.02 \\
Feature-of & 72 & .21 & .24 & $-$.03 \\
Part-of & 58 & .19 & .22 & $-$.03 \\
Evaluate-for & 43 & .16 & .18 & $-$.02 \\
Hyponym-of & 34 & .14 & .19 & $-$.05 \\
Conjunction & 21 & .11 & .08 & +.03 \\
Compare & 13 & .09 & .12 & $-$.03 \\
\bottomrule
\end{tabular}
\caption{Per-relation-type $\rho_{\text{SJ}}$ on SciERC RE ($N{=}8$). All types exhibit weak correlation. The bootstrap 95\% CI for the overall SJ$-$FK difference spans zero: [$-$0.08, +0.05].}
\label{tab:pertype}
\end{table}

A potential explanation for RE's weak consistency signal is that specific hard relation types drag down the average.
Table~\ref{tab:pertype} rules this out: \emph{all} relation types show weak $\rho$ (range: 0.09--0.28), and the SJ--FK difference is inconsistent in direction across types.
The bootstrap 95\% CI for the aggregate SJ$-$FK difference spans zero ([$-$0.08, +0.05]), meaning we cannot even establish whether SJ is better than FK for RE.

The failure is structural rather than type-specific.
Each relation triplet (head, tail, type) has three independently varying components; a single component error (e.g., wrong tail boundary) can produce a completely non-matching triplet, amplifying noise in both SJ and FK scores.
This \emph{compound error propagation} is absent in NER, where each entity is evaluated independently.

RE at $N{=}16$ (Table~\ref{tab:re_n16}) confirms the negative finding with higher statistical power: $\rho_{\text{SJ}}{=}0.231$ (bootstrap 95\% CI: [0.18, 0.28]).

\begin{table}[t]
\centering
\small
\begin{tabular}{lcccccc}
\toprule
$N$ & $\rho_{\text{SJ}}$ & $\rho_{\text{FK}}$ & $\rho_{\text{SJ}}^{c}$ & $\rho_{\text{FK}}^{c}$ & Headroom \\
\midrule
8 & .244 & .278 & .316 & .307 & 16.7pp \\
16 & .231 & .268 & .270 & .290 & 13.7pp \\
\bottomrule
\end{tabular}
\caption{RE $N$-scaling on SciERC. Superscript $c$ denotes conditional (non-zero instances only). Signal degrades with more samples.}
\label{tab:re_n16}
\end{table}

\subsection{Model-Specific Confounds}

\begin{table}[t]
\centering
\small
\begin{tabular}{lccccc}
\toprule
\textbf{Model} & \textbf{F1} & $\rho_{\text{SJ}}$ & $\rho_{\text{FK}}$ & $\Delta$ & \textbf{Parse Fail} \\
\midrule
Qwen3-8B & .667 & .383 & .259 & +.124 & 0 \\
LLaMA-3.1-8B & .591 & .280 & .310 & $-$.030 & 12 \\
\bottomrule
\end{tabular}
\caption{Model comparison on SciERC NER ($N{=}8$). SJ's advantage vanishes for LLaMA, where format instability and boundary wobble confound structural comparison.}
\label{tab:multimodel}
\end{table}

Table~\ref{tab:multimodel} tests whether SJ's advantage is model-agnostic.
On LLaMA-3.1-8B with the same dataset (SciERC NER), FK outperforms SJ ($\rho{=}0.31$ vs.\ 0.28), reversing the Qwen3-8B result.
Two factors contribute:

\paragraph{Format instability.}
LLaMA produces 12 parse failures across stochastic samples (vs.\ 0 for Qwen3-8B).
SJ interprets format inconsistency as low confidence, but these failures reflect generation style differences, not genuine uncertainty.

\paragraph{Boundary wobble.}
LLaMA's entity span boundaries exhibit more variation across samples (wider logit distribution at boundary tokens).
SJ's soft matching partially credits boundary-shifted spans, but this variation is noise---it does not correlate with accuracy.
FK's binary token-level agreement is less sensitive to boundary shifts, giving it an advantage.

These model-specific confounds mean that deploying SJ requires validation not just per dataset but per model, severely limiting its practical utility.

%------------------------------------------------------------
\section{Discussion}
\label{sec:discussion}
%------------------------------------------------------------

\subsection{Practitioner Guidelines}

Our results suggest the following conditions must \emph{all} hold for consistency-based confidence to be potentially useful for structured prediction:

\begin{enumerate}
\item \textbf{Mid-range base performance}: Greedy F1 should be in approximately [0.50, 0.80]. Below 0.50, floor effects dominate; above 0.80, ceiling effects compress the signal.
\item \textbf{Simple output structure}: NER-like tasks where errors are independent per entity. Compound structures (RE triplets) introduce error propagation that destroys the consistency--accuracy link.
\item \textbf{Stable generation format}: The model must produce well-formed outputs consistently. Format variation and boundary wobble confound structural metrics.
\item \textbf{In-distribution data}: Distribution shift causes the consistency signal to collapse ($\rho{=}0.12$ OOD vs.\ 0.19 in-domain).
\end{enumerate}

These conditions are restrictive. In practice, the scenarios where consistency is most \emph{needed} (hard instances, complex tasks, OOD data) are precisely the scenarios where it fails.

\subsection{Why Structured Prediction Is Different from Reasoning}

Self-consistency succeeds for reasoning tasks because: (a) the output space is discrete and small (a number, a multiple-choice letter), making majority voting straightforward; (b) correct reasoning paths converge deterministically on the correct answer; (c) errors are typically ``randomly wrong'' rather than ``consistently wrong.''

Structured prediction violates all three properties: (a) the output space is combinatorial (exponential in the number of entities/relations); (b) multiple structurally different outputs can be partially correct; (c) systematic model weaknesses produce consistent errors (e.g., always missing a specific entity type).

\subsection{Limitations}

Our study has several limitations.
First, we evaluate only 8B-parameter models; larger models with stronger instruction-following may exhibit different sampling behavior.
Second, our SJ metric uses a specific soft matching formulation; alternative structural similarity metrics (e.g., graph edit distance for RE) might perform differently.
Third, we evaluate on English IE datasets only; cross-lingual generalization is unknown.
Fourth, the capability stratification analysis (Section~\ref{sec:stratification}) lacks statistical power due to small per-stratum sample sizes.

%------------------------------------------------------------
\section{Conclusion}
%------------------------------------------------------------

We presented a systematic evaluation of self-consistency as a confidence signal for structured information extraction.
Our findings are predominantly negative: the consistency--accuracy correlation is fragile, dataset-specific, model-specific, and does not enable practical test-time scaling through best-of-$N$ selection.
Through error decomposition and capability stratification, we identified three failure modes---performance ceiling effects, capability floor effects, and model-specific sampling confounds---that together predict a narrow operating range for consistency-based confidence.
We hope these results help practitioners make informed decisions about when (not) to invest in consistency-based confidence estimation, and motivate future work on more robust confidence signals for structured prediction.

\begin{thebibliography}{10}

\bibitem{wang2023selfconsistency}
Xuezhi Wang, Jason Wei, Dale Schuurmans, Quoc Le, Ed Chi, Sharan Narang, Aakanksha Chowdhery, and Denny Zhou.
Self-consistency improves chain of thought reasoning in language models.
In \emph{ICLR}, 2023.

\bibitem{brown2024large}
Bradley Brown, Jordan Juravsky, Ryan Ehrlich, Ronald Clark, Quoc V. Le, Christopher R\'{e}, and Azalia Mirhoseini.
Large language monkeys: Scaling inference compute with repeated sampling.
\emph{arXiv preprint arXiv:2407.21787}, 2024.

\bibitem{chen2023codet}
Bei Chen, Fengji Zhang, Anh Nguyen, Daoguang Zan, Zhiruo Lin, Jian-Guang Lou, and Weizhu Chen.
CodeT: Code generation with generated tests.
In \emph{ICLR}, 2023.

\bibitem{kadavath2022language}
Saurav Kadavath, Tom Conerly, Amanda Askell, et al.
Language models (mostly) know what they know.
\emph{arXiv preprint arXiv:2207.05221}, 2022.

\bibitem{li2023evaluating}
Peng Li, Tianxiang Sun, Qiong Tang, Hang Yan, Yuanbin Wu, Xuanjing Huang, and Xipeng Qiu.
Evaluating ChatGPT's information extraction capabilities.
\emph{arXiv preprint arXiv:2304.11633}, 2023.

\bibitem{wadhwa2023revisiting}
Somin Wadhwa, Silvio Amir, and Byron C. Wallace.
Revisiting relation extraction in the era of large language models.
In \emph{ACL}, 2023.

\bibitem{rogers2020primer}
Anna Rogers.
A primer on negative results in NLP.
In \emph{ACL}, 2020.

\bibitem{luan2018multi}
Yi Luan, Luheng He, Mari Ostendorf, and Hannaneh Hajishirzi.
Multi-task identification of entities, relations, and coreference for scientific knowledge graph construction.
In \emph{EMNLP}, 2018.

\bibitem{tjong2003introduction}
Erik F. Tjong Kim Sang and Fien De Meulder.
Introduction to the {CoNLL}-2003 shared task: Language-independent named entity recognition.
In \emph{CoNLL}, 2003.

\bibitem{derczynski2017results}
Leon Derczynski, Eric Nichols, Marieke van Erp, and Nut Limsopatham.
Results of the {WNUT}2017 shared task on novel and emerging entity recognition.
In \emph{W-NUT}, 2017.

\end{thebibliography}

\end{document}
